\section{Pasar Finansial}
"Uang" merupakan entitas fisik yang sudah menjadi alat untuk merepresentasikan kondisi interaksi sosial bagi seluruh manusia di dunia, di mana transformasinya menjadi instrumen digital memperluas spektrum pertukaran nilai ekonomi secara global \cite{hidalgo_quantum_2006}. Evolusi mekanisme pertukaran ini membentuk sebuah arsitektur jaringan makroekonomi yang luas untuk mengatur distribusi sumber daya secara optimal melalui korelasi antar aset yang bersifat non-linear \cite{mantegna_introduction_2000}. Seluruh kumpulan dari berbagai institusi perbankan, instrumen investasi, serta mekanisme regulasi formal yang mengatur sirkulasi dan akumulasi dari arus uang tersebut disebut sebagai sistem finansial. Struktur organisasi ini memungkinkan adanya ketergantungan lintas aset yang teratur, di mana nilai dari setiap instrumen sepenuhnya ditentukan oleh konsensus informasi serta dinamika jaringan global yang terstruktur \cite{gong_quantum_2025}. Objek finansial tersebut selanjutnya dapat diperjual-belikan, dipertukarkan, dan dispekulasikan secara dinamis di dalam sebuah lingkungan interaksi ekonomi yang terminologinya disebut sebagai pasar finansial.

\subsection{Sejarah Uang}
Evolusi sistem pertukaran dalam sejarah manusia dimulai dari kebutuhan biologis yang terbatas namun berhadapan dengan keinginan yang bersifat tidak terbatas, sehingga menciptakan konflik kelangkaan yang mendasari munculnya instrumen yang mengontrol ekonomi. Mekanisme barter tersebut menyebabkan inefisiensi masif sehingga mendorong penggunaan uang komoditas serta logam mulia sebagai medium transisi yang lebih stabil, meskipun adopsi tersebut masih terhambat oleh beban logistik dan keterbatasan mobilitas fisik yang signifikan \. Pada akhirnya, manusia mulai melakukan transisi sistem transaksi menuju uang \textit{fiat} di mana nilai nominal sepenuhnya bersandar pada kepercayaan kolektif terhadap stabilitas pemerintahan. Dalam kerangka teori informasi, uang dipandang sebagai instrumen kompresi data yang sangat efisien untuk menyederhanakan rasio harga barter yang kompleks menjadi representasi numerik yang tunggal \cite{kolmogorov_three_1968}. Sejarah panjang dari peradaban manusia menunjukkan bahwa sistem moneternya selalu mengalami evolusi seiring dengan perkembangan teknologi dan struktur sosial, dimulai dari komoditas fisik, logam mulia, uang kertas, hingga kini dalam bentuk digital yang terintegrasi penuh dengan teknologi finansial (\textit{fintech}) dan sistem keuangan global berbasis jaringan komputer yang kompleks.

Dalam ekosistem finansial, uang harus dipahami sebagai alat tukar yang aktif dan bukan sekadar instrumen penyimpan kekayaan yang bersifat statis serta tidak produktif. Inflasi muncul sebagai fenomena makroekonomi yang ditandai dengan kenaikan tingkat harga secara berkelanjutan. Inflasi dapat mengakibatkan degradasi daya beli serta mengikis nilai riil dari mata uang yang beredar di pasar secara signifikan. Kondisi makroekonomi ini mencerminkan adanya luapan likuiditas yang berlebihan sehingga menciptakan ketidakseimbangan antara suplai dan permintaan \cite{saslow_economic_1999}. Untuk mendeteksi tingkat inflasi, investor menggunakan parameter volatilitas inflasi sebagai indikator risiko yang sangat krusial karena sifatnya yang stokastik dan sering kali memiliki korelasi yang signifikan terhadap ekspektasi imbal hasil aset di pasar modal \cite{schwert_chapter_2003}. Parameter volatilitas inflasi ini dipengaruhi oleh faktor eksternal seperti geopolitik, bencana alam, dan kebijakan pemerintah, dapat mempengaruhi ekspektasi investor terhadap imbal hasil aset di pasar modal.

Dampak inflasi dapat merambat ke seluruh sektor investasi karna terdapat mekanisme distorsi pada sinyal harga yang dapat mengganggu ekspektasi imbal hasil riil serta memaksa agen ekonomi untuk melakukan realokasi modal setiap periode waktu \cite{schwert_chapter_2003}. Kekayaan riil seorang investor akan mengalami degradasi secara eksponensial apabila laju inflasi tahunan melebihi tingkat suku bunga tabungan nominal, sehingga menciptakan insentif kuat bagi pelaku pasar untuk beralih dari aset kas menuju instrumen yang memiliki produktivitas lebih tinggi. Fenomena ini sering kali dipandang sebagai pajak tersembunyi yang mendorong pergerakan modal menuju aset berisiko \cite{saslow_economic_1999}. Respon asimetris investor terhadap potensi kerugian riil tersebut dapat dijelaskan melalui kerangka perilaku dalam teori prospek \cite{kahneman_prospect_1979}. Ketidakpastian inflasi yang tinggi pada akhirnya menuntut implementasi strategi diversifikasi portofolio untuk memitigasi risiko penurunan daya beli sekaligus mempertahankan stabilitas nilai kekayaan melalui pemilihan kombinasi aset yang efisien \cite{markowitz_portfolio_1952, demiguel_optimal_2009}.

\subsection{Investasi}
Strategi diversifikasi pada dasarnya diimplementasikan untuk melindungi kekayaan dari dampak inflasi dengan menginvestasikan beberapa aset dengan masing-masing alokasi yang terukur dari total modal yang ingin diinvestasikan. Investasi didefinisikan sebagai tindakan strategis dalam mengeluarkan sumber daya finansial pada saat ini dengan ekspektasi kuat untuk memperoleh imbal hasil atau apresiasi nilai yang signifikan di masa depan \cite{abdul_hali_markowitz_2020}. Aktivitas ekonomi ini melibatkan pengorbanan konsumsi jangka pendek untuk mengakumulasi modal secara produktif untuk memenuhi kebutuhan jangka panjang bagi para agen ekonomi yang bertindak secara rasional. Karena urgensi investasi menjadi semakin nyata di tengah sistem moneter yang mengalami tekanan inflasi persisten, maka penyimpanan kas secara pasif akan mengakibatkan degradasi kekayaan riil secara berkelanjutan \cite{schwert_chapter_2003}. Dengan pemilihan instrumen yang tepat dan terukur, agen ekonomi dapat memitigasi risiko penurunan daya beli sekaligus berpartisipasi aktif dalam pertumbuhan produktivitas sektor riil maupun dinamika pasar finansial global \cite{markowitz_portfolio_1952}. Oleh karena itu, investasi yang terencana dengan baik berfungsi sebagai mekanisme pertahanan utama bagi individu maupun institusi dalam menghadapi ketidakpastian ekonomi serta fluktuasi pasar yang sering kali bersifat stokastik dan sulit diprediksi secara linear.

Spektrum strategi investasi sangat bervariasi, mulai dari pendekatan konvensional yang mengutamakan preservasi modal hingga strategi agresif yang mengejar pertumbuhan nilai aset secara eksponensial dalam rentang waktu yang relatif singkat. Keragaman strategi ini berlandaskan pada prinsip risk-return trade-off, yang menyatakan bahwa ekspektasi imbal hasil yang tinggi selalu berbanding lurus dengan tingkat risiko atau ketidakpastian sistemik yang harus ditanggung \cite{markowitz_portfolio_1952}. Investor, sebagai agen rasional, berusaha mencapai utilitas maksimal dengan memilih berbagai instrumen finansial yang sesuai dengan batas toleransi risiko untuk mengoptimalkan performa portofolio melalui pendekatan alokasi aset yang efisien \cite{sikalo_efficient_2022}. Implementasi strategi diversifikasi menjadi instrumen utama dalam menyeimbangkan proporsi antara aset berisiko dan aset aman untuk meminimalkan dampak volatilitas kolektif sekaligus mempertahankan efisiensi operasional modal secara berkelanjutan \cite{demiguel_optimal_2009}. Dalam pemodelan ekonomi modern, parameter penghindaran risiko (\textit{risk aversion}) bersifat endogen karena nilainya terus berevolusi mengikuti fluktuasi kondisi kekayaan serta pengalaman historis yang dialami oleh para agen di pasar modal secara dinamis \cite{levy_arrow-pratt_1991}. Landasan perilaku ini didukung oleh Teori Prospek yang menjelaskan bahwa para agen cenderung menunjukkan respon psikologis yang asimetris, di mana rasa sakit akibat kerugian dirasakan jauh lebih besar dibandingkan kepuasan yang diperoleh dari keuntungan investasi \cite{kahneman_prospect_1979}. Walaupun demikian, evolusi parameter risiko dalam ekosistem pasar dapat memicu terjadinya transisi fase yang dramatis mulai dari kondisi likuiditas tinggi hingga mencapai titik stagnasi atau pembekuan aktivitas ekonomi secara total \cite{galam_ising_2010}. Aspek psikologis dan perubahan akumulasi kekayaan diinegrasikan ke dalam parameter risiko dengan tujuan agar strategi investasi yang dibentuk lebih realistis dalam mendeteksi potensi krisis finansial melalui identifikasi sinyal peringatan dini \cite{gong_quantum_2025}. 